% ============================================================
%  MCM/ICM (美赛) 论文 LaTeX 英文模板
%  2026 官方规则要点：
%  - 全文（含 Summary Sheet / 目录 / 参考文献 / 附录 / 代码）≤ 25 页硬上限
%  - 字号 ≥ 12pt；每页页眉 "Team #XXXXX, Page X of 25"
%  - Summary Sheet 为第 1 页，约 300 词，占页 1/2-2/3，严禁溢出
%  - 附录代码计入 25 页 → 只放核心脚本，其余移入支撑材料
%  - AI 使用报告为单独文件（不计入 25 页），须如实披露 AI 使用
%  编译：xelatex main.tex（参考文献为内联 thebibliography，无需 bibtex；
%        若改用 .bib + \bibliography 管线，才需要 bibtex 两轮重编译）
% ============================================================
\documentclass[12pt,letterpaper]{article}

% ---------- 语言与字体 ----------
\usepackage[utf8]{inputenc}
\usepackage[T1]{fontenc}
% ⚠️ 加载顺序铁律（2026 实测验证）：
%   ① amsmath/amsthm 必须【先于】newtxmath——否则 \openbox 冲突
%     （Command `\openbox' already defined）
%   ② 加载 newtxmath 后【不得再加载 amssymb】——两者都定义 \Bbbk
%     （newtxmath 已内置全部 AMS 符号）
\usepackage{amsmath, amsthm, mathtools}
\usepackage{newtxtext, newtxmath}  % Times 系字体（12pt 可读性优先）

% ---------- 页面设置（US Letter；官方建议：上/下 2.5cm，左/右 3cm）----------
\usepackage{geometry}
\geometry{letterpaper, top=2.5cm, bottom=2.5cm, left=3cm, right=3cm}

% ---------- 页眉：Team #XXXXX, Page X of N ----------
% ⚠️ v2 修正：总页数用 \pageref{LastPage} 动态取（lastpage 宏包），
%    旧版硬编码 "of 25"——实际页数少于 25 时页眉与真实总页数不符
\usepackage{fancyhdr}
\usepackage{lastpage}
\pagestyle{fancy}
\fancyhf{}
\fancyhead[L]{\small Team \#<TEAM\_ID>}          % Replace <TEAM\_ID> before delivery.
\fancyhead[R]{\small Page \thepage\ of \pageref{LastPage}}
\renewcommand{\headrulewidth}{0.4pt}

% ---------- 行距 ----------
\usepackage{setspace}
\setstretch{1.15}

% （数学宏包已在字体区按顺序加载：amsmath/amsthm → newtx，见上方注释）

% ---------- 表格 ----------
\usepackage{booktabs, longtable, multirow, multicol, array}
\usepackage[table]{xcolor}

% ---------- 图片 ----------
\usepackage{graphicx}
\graphicspath{{figures/}{./}}
\usepackage{float}
\usepackage{caption, subcaption}
\captionsetup{font={small}, labelfont=bf}

% ---------- 代码块（附录代码计入 25 页 → 精简）----------
\usepackage{listings}
\lstset{
    basicstyle=\small\ttfamily,
    keywordstyle=\color{blue},
    commentstyle=\color{gray},
    stringstyle=\color{orange},
    numbers=left,
    numberstyle=\tiny\color{gray},
    frame=single,
    breaklines=true,
    tabsize=4,
    language=Python
}

% ---------- 算法 ----------
\usepackage{algorithm, algpseudocode}

% ---------- 超链接 ----------
\usepackage{hyperref}
\hypersetup{colorlinks=true, linkcolor=black, citecolor=black, urlcolor=blue}

\begin{document}

% ============================================================
%  SUMMARY SHEET (Page 1 of 25) —— 官方建议约 300 词，占页 1/2-2/3
% ============================================================
\begin{center}
  {\Large\bfseries 2026 MCM/ICM Summary Sheet}\\[6pt]
  {\large Problem Chosen: \underline{\ \ <PROBLEM\_ID>\ \ } \quad
   Team Control Number: \underline{\ \ <TEAM\_ID>\ \ }}\\[10pt]
  {\LARGE\bfseries <PAPER\_TITLE>}\\[4pt]
  {\normalsize (Team \#<TEAM\_ID>)}
\end{center}

% 摘要正文：背景(1-2句) → 核心建模方法(2句) → 主要发现(2-3句) → 模型优势(1句)
% 必须含每问的关键数值 + 至少 1 问交叉验证；约 250-300 词
In <PROBLEM\_BACKGROUND>, this paper develops <CORE\_MODELING\_FRAMEWORK>.

For Problem 1, we propose <METHOD\_NAME> based on <THEORY>, obtaining
<KEY\_RESULT\_WITH\_UNIT> with 95\% CI <CONFIDENCE\_INTERVAL>, validated against
<ALTERNATIVE\_METHOD> with relative error <RELATIVE\_ERROR>.

For Problem 2, <PROBLEM\_2\_RESULT\_AND\_LIMITATION>.

The main contributions are: (1) <CONTRIBUTION\_1>; (2) <CONTRIBUTION\_2>;
(3) <CONTRIBUTION\_3>. Sensitivity analysis indicates that the conclusions
are robust within <SUPPORTED\_PARAMETER\_RANGE>.

\vspace{0.5em}
\noindent\textbf{Keywords:} <KEYWORD\_1>; <KEYWORD\_2>; <KEYWORD\_3>

\newpage

% ============================================================
%  MAIN BODY —— ⛔ 全文（含本页至附录结束）合计 ≤ 25 页
% ============================================================

% ---------- 1. Problem Restatement ----------
\section{Problem Restatement}
\subsection{Background}
\subsection{Problem Decomposition}
% 逐问重述 + 数学化表述 + 已有方法局限

% ---------- 2. Assumptions ----------
\section{Model Assumptions}
\begin{enumerate}
  \item Assumption 1: ... (justification)
  \item Assumption 2: ... (justification)
\end{enumerate}

% ---------- 3. Notation ----------
\section{Notation}
\begin{table}[H]
  \centering
  \caption{Notation}\label{tab:notation}
  \begin{tabular}{clc}
    \toprule
    \textbf{Symbol} & \textbf{Meaning} & \textbf{Unit} \\
    \midrule
    $x$ & ... & ... \\
    \bottomrule
  \end{tabular}
\end{table}

% ---------- 4. Models ----------
\section{Model Formulation and Solution}
\subsection{Problem 1: <PROBLEM\_1\_TITLE>}
\subsubsection{Model Derivation: <MODEL\_DERIVATION\_TITLE>}
\subsubsection{Theoretical Analysis: <THEORETICAL\_ANALYSIS\_TITLE>}
\subsubsection{Algorithm Design: <ALGORITHM\_DESIGN\_TITLE>}
\begin{algorithm}[H]
\caption{<ALGORITHM\_CAPTION>}\label{alg:<ALGORITHM\_LABEL>}
\begin{algorithmic}[1]
\Require <ALGORITHM\_INPUT>
\Ensure <ALGORITHM\_OUTPUT>
\State <ALGORITHM\_STEP\_1>
\State <ALGORITHM\_STEP\_2>
\end{algorithmic}
\end{algorithm}
\subsubsection{Ablation Study: <ABLATION\_TITLE>}
\subsubsection{Results and Uncertainty Quantification: <RESULTS\_TITLE>}

\subsection{Problem 2: <PROBLEM\_2\_TITLE>}
% Reuse the same evidence-backed structure only when it applies; otherwise
% state the justified NOT_APPLICABLE decision and its alternative evidence.

% ---------- 5. Validation & Sensitivity (美赛强制章节) ----------
\section{Model Validation and Sensitivity Analysis}
\subsection{Cross-Validation}
\subsection{Sensitivity Analysis}
\subsection{Robustness (Monte Carlo)}

% ---------- 6. Strengths & Weaknesses (美赛强制章节) ----------
\section{Strengths and Weaknesses}
\subsection{Strengths}
\subsection{Weaknesses}

% ---------- 7. Discussion ----------
\section{Discussion}
\subsection{Comparison with Literature}
\subsection{Unexpected Findings}
\subsection{Limitations and Future Work}

% ---------- 8. Conclusion ----------
\section{Conclusion}
\label{lastmain}  % ⛔ 与锁14同构：正文末页标记（美赛按"正文+附录≤25页"判定）
% 核心贡献 + 数值汇总 + 决策建议

% ---------- References (IEEE/APA) ----------
% 内联 thebibliography：无需 bibtex；如切换为 .bib 管线，
% 再取消下一行注释并删除 thebibliography 环境。
% \bibliographystyle{IEEEtran}   % 仅在 \bibliography{refs} 管线下生效
\begin{thebibliography}{99}
\bibitem{ref1} Author. Title[J]. Journal, Year, Vol(Issue): Pages.
\end{thebibliography}

% ============================================================
%  APPENDIX —— ⛔ 计入 25 页上限：只放核心脚本，其余移入支撑材料
%  并在附录开头注明 "Full code: see supporting materials (zip)"
% ============================================================
\newpage
\appendix
\section*{Appendix: Core Code}
% [核心脚本精简版；注释写明"对应论文 §X.X 节"与"对应公式(N)"]
% 全量代码 → 支撑材料 zip（文件名与论文附录文件列表一致）

\end{document}
